\section{Triangular selection in the probe limit}
\label{app:probe}

We prove the statement used in section~\ref{sec:5.2}: in the probe limit,
near the onset, the triangular lattice is the global minimum of the quartic
free energy over all one-flux-quantum Bravais lattices. Once the kernel is
shown to be completely monotone this is a corollary of Montgomery's
theorem; what is specific to holography is the Stieltjes representation of
the kernel below. The question was left open
in~\cite{Bu:2012mq}, which compared rhombic candidates numerically.

\subsection{The reduced free energy}
\label{app:probe-free-energy}

On $\mathrm{AdS}_5$--Schwarzschild in the
coordinate $u = 1/r$, with $f = 1 - u^4$ and $u_H = 1$, the static planar
ansatz $W_x = \eta\,w_0(u)\psi(x,y)$, $W_y = -iW_x$,
$A^3 = (Bx + \eta^2 a_y)\,dy + \eta^2 a_x\,dx$ in the gauge
$W_u = 0$ reduces the Yang--Mills free energy density exactly, to all orders
in the fields, to
\begin{equation}
\mathcal{F} = \frac{1}{2u}\Big[|F^W_{xy}|^2 + \big(B + \hat b + \mathrm{Im}(\bar W_xW_y)\big)^2\Big] +
\frac{f}{2u}\Big[|W_x'|^2 + |W_y'|^2 + a_x'^2 + a_y'^2\Big]~,
\end{equation}
with $\hat b = \partial_xa_y - \partial_ya_x$ the induced field and $F^W_{xy} = D_xW_y - D_yW_x$
covariant under the full $A^3$. Two facts make this tractable. On the
polarised lowest Landau level $F^W_{xy}$ vanishes at $a = 0$ and is
$O(\eta^3)$ with the induced field, so the condensate pays a gradient
cost in $u$ only; and $\mathrm{Im}(\bar W_xW_y) = -|W_x|^2$, so the
condensate is diamagnetic at contact level and the local field is
$B + \eta^2(\hat b - w_0^2|\psi|^2)$. The $u$-components of the Yang--Mills
equations close on the ansatz. The expansion in $\eta$ then has no odd
terms. Its quadratic term is the Sturm--Liouville form of
section~\ref{sec:3} and is exactly linear in $B$ at fixed $w_0$,
$F_2 = -\eta^2(B - B_c)J_2$ with $J_2 = \int w_0^2\,du/u = 1$, and its
quartic term is
\begin{equation}
\mathcal{F}_4 = \frac{\big(\hat b - w_0^2|\psi|^2\big)^2}{2u} + \frac{f}{2u}\big(a_x'^2 + a_y'^2\big)~.
\end{equation}
Expanding $|\psi|^2 = \sum_G\omega_Ge^{iG\cdot x}$ on the reciprocal
lattice, each harmonic decouples. The induced field at harmonic $G$ obeys
$-(f\hat b_G'/u)' + (G^2/u)\,\hat b_G = (G^2/u)\,\omega_Gw_0^2$ with $\hat b_G(0) = 0$ and
regularity at the horizon, and the on-shell exchange energy is minus half
the source coupling. This is the abelian sector of appendix~\ref{app:response-abelian} with
$p_1 = 1/u$ and $p_2 = f/u$, and it gives the probe kernel
$K(s;0) = C_4 - X(s)$ with $s = G^2$ (in the probe limit the coupled and abelian-only screenings $X$ and $X_{\rm ab}$ of section~\ref{sec:5.1} coincide), $C_4 = \int w_0^4\,du/u$ and
$X(s) = \int \hat b_G w_0^2\,du/u$ at unit source. The operator $\mathcal{T}$ of
appendix~\ref{app:response-abelian} is here the onset operator of section~\ref{sec:3} itself, so its
spectrum is the tower $B_n$ of radial overtones of the linear problem, all in
the lowest Landau level, and the spectral
form reads
\begin{equation}
C_4 - X(s) = \sum_n c_n^2\,\frac{B_n}{B_n + s}~, \qquad c_n = \langle w_0^2,\varphi_n\rangle~, \qquad \sum_n c_n^2 = C_4~.
\end{equation}
Hence $0 < X(s) < C_4$, $X(0) = 0$, $X \to C_4$ as $s \to \infty$, and
$(-1)^k\,d^k(C_4 - X)/ds^k \ge 0$ for every $k$: the screened repulsion is
completely monotone. The $G = 0$ harmonic of $\hat b$ vanishes by the Bianchi
identity at fixed applied field, so $K_{G=0} = C_4$ in the probe limit.

\subsection{The lowest-Landau-level identity}
\label{app:probe-lll-identity}

For an element of the lowest
Landau level with one flux quantum per unit cell, $|\omega_G| = e^{-G^2/4B}$
for every Bravais lattice, which is the origin of the weights
$e^{-\gamma_G/2}$ in section~\ref{sec:5.1}. This is the classical result of
the flux-line-lattice literature, stated for general lattice symmetry as
equation (1.6) of~\cite{Brandt:1995rpp} and equation (B6)
of~\cite{Brandt:2003prb}, where the coefficients are
$(-1)^{m+mn+n}e^{-G_{mn}^2 S/8\pi}$ with $S = 2\pi/B$ the cell area. The
lowest-Landau-level projection of $e^{iG\cdot x}$ carries the Gaussian form
factor. With one flux quantum per cell the remaining guiding-centre factor
is a magnetic translation by a lattice vector, of which the state is an
eigenstate with eigenvalue of unit modulus. The identity reproduces the
Abrikosov ratios $\beta_\square = 1.1803406$ and
$\beta_\triangle = 1.1595953$ (the $\beta_A$ of the vortex literature,
unrelated to our backreaction parameter), and we checked it numerically on
a generic lattice as well. The quartic functional is therefore the one of section~\ref{sec:5.1} at
$\beta = 0$,
\begin{equation}
F = -\eta^2(B - B_c) + \Big[\tfrac12C_4 + C(\tau)\Big]\eta^4~, \qquad
C(\tau) = \tfrac12\sum_{G\in\Lambda^*(\tau)\setminus0} e^{-G^2/2B_c}\,\big(C_4 - X(G^2)\big)~,
\end{equation}
minimised over the amplitude by $\eta_\star^2 = (B - B_c)/(C_4 + 2C)$
with $F_{\min} = -(B - B_c)^2/(2C_4 + 4C)$, so the transition is second
order and the lattice that forms is the one that minimises $C(\tau)$.

\subsection{The proposition}
\label{app:probe-theorem}

Since $C_4 - X$ is a Stieltjes transform of a positive
measure, the weighted kernel is a Laplace transform of one,
\begin{equation}
k(s) \equiv e^{-s/2B_c}\,\big(C_4 - X(s)\big) = \int_{1/2B_c}^\infty e^{-st}\,d\mu(t)~, \qquad
d\mu(t) = \sum_n c_n^2B_n\,e^{-B_n(t - 1/2B_c)}\,dt~,
\end{equation}
and the lattice-dependent part of the quartic coefficient is a positive
superposition of Gaussian lattice theta functions,
\begin{equation}
C(\tau) = \tfrac12\sum_{G\neq0}k(|G|^2) = \tfrac12\int d\mu(t)\,\big[\Theta_{\Lambda^*}(t) - 1\big]~, \qquad
\Theta_\Lambda(t) = \sum_{v\in\Lambda}e^{-t|v|^2}~.
\end{equation}
Flux quantisation fixes the covolume of $\Lambda^*$ at $2\pi B_c$, so
$\Lambda^*$ ranges over all two-dimensional lattices of fixed covolume, and
the triangular lattice is its own dual shape. Montgomery's
theorem~\cite{Montgomery:1988mt} states that among these the Gaussian theta
function is minimised, for every $t > 0$, by the triangular lattice alone;
integrating against the positive measure $\mu$ gives
$C(\tau) > C(\tau_\triangle)$ for every non-triangular $\tau$. That a
completely monotone interaction is minimised by the triangular lattice in
this way was observed by Cohn and Kumar~\cite{Cohn:2007universally} and
is Proposition 3.1 of~\cite{Betermin:2016theta}; Proposition 3.4 there
gives a convex decreasing positive interaction for which the square lattice
has lower energy than the triangular one at some densities, which is why complete monotonicity, and not
positivity, is the property the argument needs.

The proof used three things only: the identity of appendix~\ref{app:probe-lll-identity}, the positivity and
self-adjointness of the radial operator, and Montgomery's theorem. The
numbers $C_4$, $c_n$ and $B_n$, and with them the temperature and the
blackening function, never entered, so the result holds at every
temperature near $B_c$ and on any background that preserves the structure
of the reduced functional. That is why it goes through unchanged for the
abelian sector on the brane, appendix~\ref{app:response-abelian}, and why it is no longer
guaranteed once the metric sector enters, section~\ref{sec:5.2}. Its scope is the class of
Abrikosov's analysis: Bravais lattices with one flux quantum per cell.
Multi-quantum cells and non-lattice lowest-Landau-level configurations are
not excluded by it, although the triangular lattice is locally stable
against all lowest-Landau-level perturbations at one flux quantum
(section~\ref{sec:5.3}). For the flat-space lowest-Landau-level functional,
whose Abrikosov ratio is a Gaussian theta function, the same route to the
triangular lattice is known~\cite{Nier:2007theta,Aftalion:2006lll}, with
the Abrikosov case proved directly in~\cite{Nonnenmacher:1998chaotic};
what is new here is that the holographic kernel keeps the structure.

\subsection{Numbers}
\label{app:probe-numbers}

The values below use $u_H = 1$ and $J_2 = 1$. The brane normalisation
$\int p_1 w_0^2\,dr = 1$ of section~\ref{sec:5} reduces to $J_2 = 1$ at
$\beta = 0$ and differs from it by a positive factor at $\beta > 0$, which is
why the brane results of section~\ref{sec:5.2} are quoted as ratios.

\begin{center}
\begin{tabular}{lc}
\toprule
quantity & value \\
\midrule
$C_4$ & $1.5805931$ \\
$C(\tau_\triangle)$ & $0.0184443$ \\
$C(\tau_\square)$ & $0.0229818$ \\
$(C_\square - C_\triangle)/C_\triangle$ & $+0.2460$ \\
\bottomrule
\end{tabular}
\end{center}
The relative margin $+0.246$ is the probe value against which the brane
margins of section~\ref{sec:5.2} are compared; it is essentially unchanged by the
abelian sector of the throat ($+0.250$) and reduced by the metric sector
($+0.1766$). A grid scan of $C(\tau)$ over the fundamental domain
up to $\tau_2 = 5$, the $\beta = 0$ row of the moduli scan of
appendix~\ref{app:numerics-scan}, finds its minimum at the
triangular point, which checks the proposition for this kernel.
